\documentclass[10pt,a4paper]{amsart}
\usepackage{amsmath,amssymb,amsthm,graphicx}
\usepackage{verbatim}
\newenvironment{reference}{\comment}{\endcomment}
%\newenvironment{reference}{}{}
\newenvironment{slogan}{\comment}{\endcomment}
\newenvironment{history}{\comment}{\endcomment}

% For commutative diagrams we use Xy-pic
\usepackage[all]{xy}

% We use 2cell for 2-commutative diagrams.
\xyoption{2cell}
% Theorem environments.
%
\theoremstyle{plain}
\newtheorem{theorem}[subsection]{Theorem}
\newtheorem{proposition}[subsection]{Proposition}
\newtheorem{lemma}[subsection]{Lemma}

\theoremstyle{definition}
\newtheorem{definition}[subsection]{Definition}
\newtheorem{example}[subsection]{Example}
\newtheorem{exercise}[subsection]{Exercise}
\newtheorem{situation}[subsection]{Situation}

\theoremstyle{remark}
\newtheorem{remark}[subsection]{Remark}
\newtheorem{remarks}[subsection]{Remarks}

\numberwithin{equation}{subsection}

% Macros
%
\def\lim{\mathop{\mathrm{lim}}\nolimits}
\def\colim{\mathop{\mathrm{colim}}\nolimits}
\def\Spec{\mathop{\mathrm{Spec}}}
\def\Hom{\mathop{\mathrm{Hom}}\nolimits}
\def\Ext{\mathop{\mathrm{Ext}}\nolimits}
\def\SheafHom{\mathop{\mathcal{H}\!\mathit{om}}\nolimits}
\def\SheafExt{\mathop{\mathcal{E}\!\mathit{xt}}\nolimits}
\def\Sch{\mathit{Sch}}
\def\Mor{\mathop{\mathrm{Mor}}\nolimits}
\def\Ob{\mathop{\mathrm{Ob}}\nolimits}
\def\Sh{\mathop{\mathit{Sh}}\nolimits}
\def\NL{\mathop{N\!L}\nolimits}
\def\CH{\mathop{\mathrm{CH}}\nolimits}
\def\proetale{{pro\text{-}\acute{e}tale}}
\def\etale{{\acute{e}tale}}
\def\QCoh{\mathit{QCoh}}
\def\Ker{\mathop{\mathrm{Ker}}}
\def\Im{\mathop{\mathrm{Im}}}
\def\Coker{\mathop{\mathrm{Coker}}}
\def\Coim{\mathop{\mathrm{Coim}}}

% Boxtimes
%
\DeclareMathSymbol{\boxtimes}{\mathbin}{AMSa}{"02}

%
% Macros for moduli stacks/spaces
%
\def\QCohstack{\mathcal{QC}\!\mathit{oh}}
\def\Cohstack{\mathcal{C}\!\mathit{oh}}
\def\Spacesstack{\mathcal{S}\!\mathit{paces}}
\def\Quotfunctor{\mathrm{Quot}}
\def\Hilbfunctor{\mathrm{Hilb}}
\def\Curvesstack{\mathcal{C}\!\mathit{urves}}
\def\Polarizedstack{\mathcal{P}\!\mathit{olarized}}
\def\Complexesstack{\mathcal{C}\!\mathit{omplexes}}
% \Pic is the operator that assigns to X its picard group, usage \Pic(X)
% \Picardstack_{X/B} denotes the Picard stack of X over B
% \Picardfunctor_{X/B} denotes the Picard functor of X over B
\def\Pic{\mathop{\mathrm{Pic}}\nolimits}
\def\Picardstack{\mathcal{P}\!\mathit{ic}}
\def\Picardfunctor{\mathrm{Pic}}
\def\Deformationcategory{\mathcal{D}\!\mathit{ef}}
\usepackage{cleveref}
\crefname{originalreference}{original reference}{original references}
\Crefname{originalreference}{Original reference}{Original references}
\makeatletter
\newlabel{sites-lemma-localize-morphism}{{03CF}{source}{}{sourcepointer-1}{}}
\newlabel{sites-lemma-localize-morphism@cref}{{[originalreference][0][]03CF}{[1][0][]source}{}{}{}}
\makeatother
\title{Localizing a morphism of ringed sites: the canonical (f')\textasciicircum{}\textbackslash{}sharp and the commutative square of ringed topoi}
\author{The Stacks Project Authors}
\date{stacks-project:sites-modules.tex:local-snapshot-20261012; GNU FDL 1.2 or later}
\begin{document}
\maketitle
\noindent\textit{Editorial note.} All delivered mathematical text is the author's own, copied verbatim from the pinned Stacks Project chapter file sites-modules.tex; no statement or proof step was rewritten, reordered or removed.\par\medskip
\noindent\textit{Editorial note.} The counted span contains the author's phrase 'everything else follows from' with a precise citation of Sites, Lemma sites-lemma-localize-morphism; kept verbatim and recorded here as the author's own reduction step, not as an omission.\par\medskip
\noindent\textit{Editorial note.} The counted span cites one Sites lemma, recorded in assembly.external\_references with canonical tags as navigation pointers.\par\medskip
\noindent\textit{Allowed standard prerequisites (editorial).}\begin{itemize}
\item Localization square for a morphism of topoi. Stacks Project, Sites, Lemma sites-lemma-localize-morphism. Supplies the commutative square of topoi and the functor behaviour that the ringed statement refines; also underlies the identification (f')\textasciicircum{}\{-1\}j\_V\textasciicircum{}\{-1\} = j\_U\textasciicircum{}\{-1\}f\textasciicircum{}\{-1\} used in the composite.
\item j\textasciicircum{}\{-1\} = j\textasciicircum{}* on sheaves of modules for a localization morphism. Stacks Project, Modules on Sites, discussion of localization morphisms. Used in the author's closing remark to pass from the first functor equality to the second.
\end{itemize}\medskip
\ifdefined\hypertarget\hypertarget{sourcepointer-1}{}\fi
\noindent\textit{Source reference.} Original reference 03CF: sources/stacks-chow/tags.txt, source lines 4382--4382. Cross-chapter standard prerequisite cited by the author inside the counted span; it is used as a pointer and its own proof is not part of this item.\par\medskip
\begin{lemma}
\label{lemma-localize-morphism-ringed-sites}
Let
$(f, f^\sharp) :
(\mathcal{C}, \mathcal{O})
\longrightarrow
(\mathcal{D}, \mathcal{O}')$
be a morphism of ringed sites where $f$ is given by the continuous
functor $u : \mathcal{D} \to \mathcal{C}$.
Let $V$ be an object of $\mathcal{D}$ and set $U = u(V)$.
Then there is a canonical map of sheaves of rings $(f')^\sharp$
such that the diagram of
Sites, Lemma \ref{sites-lemma-localize-morphism}
is turned into a commutative diagram of ringed topoi
$$
\xymatrix{
(\Sh(\mathcal{C}/U), \mathcal{O}_U)
\ar[rr]_{(j_U, j_U^\sharp)} \ar[d]_{(f', (f')^\sharp)} & &
(\Sh(\mathcal{C}), \mathcal{O})
\ar[d]^{(f, f^\sharp)} \\
(\Sh(\mathcal{D}/V), \mathcal{O}'_V)
\ar[rr]^{(j_V, j_V^\sharp)} & &
(\Sh(\mathcal{D}), \mathcal{O}').
}
$$
Moreover, in this situation we have $f'_*j_U^{-1} = j_V^{-1}f_*$
and $f'_*j_U^* = j_V^*f_*$.
\end{lemma}

\begin{proof}
Just take $(f')^\sharp$ to be
$$
(f')^{-1}\mathcal{O}'_V =
(f')^{-1}j_V^{-1}\mathcal{O}' =
j_U^{-1}f^{-1}\mathcal{O}' \xrightarrow{j_U^{-1}f^\sharp}
j_U^{-1}\mathcal{O} = \mathcal{O}_U
$$
and everything else follows from
Sites, Lemma \ref{sites-lemma-localize-morphism}.
(Note that $j^{-1} = j^*$ on sheaves of modules if $j$ is a localization
morphism, hence the first equality of functors implies the second.)
\end{proof}

\end{document}
